% ===================== 听力 Listening =====================
\ptesection{听力}{Listening · SST + MCM + FIB + MCS + SMW + HIW + WFD}{33}

\vspace{1mm}
{\small\color{gray} 依据「听力1 / 听力2」逐题誊录，并用「听力0 / 听力详解1」补全 SST 评分明细、用长截图补全 PDF 中被悬浮按钮遮挡处。
本套听力共 \textbf{15} 题：SST 概述口语$\times$1、MCM-L 多选$\times$1、FIB-L 听力填空$\times$3、MCS-L 单选$\times$1、SMW 选词$\times$2、HIW 找错词$\times$3、WFD 听写$\times$4。}
\vspace{1mm}

\begin{notebox}
{\footnotesize 标注沿用阅读：{\color{cFix}\ding{51}\,绿=对/命中}；{\color{cWrong}\ding{55}\,红删=错}，其后 {\color{cFix}(答案/听: 正确)}。
\textbf{HIW}（找错词）：屏幕上与录音不符的词用 {\color{cWrong}\uline{红下划线}} 标出（即应点选的"错词"），括注实际读音。
\textbf{WFD}（听写）：{\color{cFix}绿}=命中词，{\color{cWrong}\sout{红删}}=多写/听错的词。}
\end{notebox}

% =================== Q1 ===================
\begin{qblock}{Q1\quad SST (C) \#21\ \ Role of Education}{概述口语 Summarize Spoken Text · 评分 \textcolor{cAccent}{\bfseries 7 / 10}}
\ptesub{你的概述 Your Summary}
\begin{promptbox}
The education is important. {\color{cAccent}\uline{The society}} has some {\color{cAccent}\uline{journal}} challenges. The influence {\color{cAccent}\uline{contain}} society and gender. The community has some critical skills.
\end{promptbox}
{\footnotesize\color{gray}（AI 标注待改：The/The society 多余冠词；journal 疑为听错词（general?）；contain → contains 主谓一致。）}
\ptesub{AI 评分明细（本题 10 分，得 \textcolor{cAccent}{\bfseries 7} 分）}
\begin{notebox}\small
\begin{tabular}{@{}lcl@{}}
Content（内容）& \textbf{2 / 2} & 内容很棒\\
Form（字数形式）& \textbf{2 / 2} & 很棒，这是必拿分数\\
Grammar（语法）& \textcolor{cWrong}{\bfseries 0 / 2} & SST 中占分比重、却很容易提高；丢分较多，建议针对性练习\\
Spelling（拼写）& \textbf{2 / 2} & 拼写很棒，再接再厉\\
Vocabulary（词汇）& \textbf{1 / 2} & PTE 不要求难词，但需用词恰当\\
\end{tabular}
\end{notebox}
{\footnotesize\color{gray}（注：SST 音频原文未在复盘截图中给出，故此处只保留你的概述与评分明细。）}
\end{qblock}

% =================== Q2 ===================
\begin{qblock}{Q2\quad MCM-L \#20\ \ Re-adjusting to Life}{多选 Multiple Answers · 评分 \textcolor{cAccent}{\bfseries 1 / 2}}
\ptesub{录音原文 Transcript}
\begin{promptbox}
Re-adjusting to life in your own country after living abroad for some time is a little like recovering from jet-lag after a long flight across several time zones. It takes time. And research indicates that after nine years living in a foreign country you never really do re-adjust. Of course, things have changed: governments have come and gone, what you knew as countryside has become a suburb, new technologies have changed the way people go about their daily lives, and so on. These changes may well have been taking place in your adopted country, but they were happening while you were there, so you could adapt as you went along. Those are not the main difficulties, however. It is in the smaller, everyday things that you might experience what is known as culture shock, although it's not really a shock, but puzzling all the same. For example, the precise way to behave at a supermarket checkout may have changed. And in ordinary conversation, the frames of reference have changed, and quite often you find that you don't really know what people are talking about, even though they are speaking your native tongue.
\end{promptbox}
\zh{在国外生活一段时间后回到自己国家重新适应，有点像长途跨时区飞行后倒时差——需要时间。研究表明，在国外生活九年后，你其实永远无法真正重新适应。当然，一切都变了：政府换了一届又一届，你记忆中的乡村变成了郊区，新技术改变了人们的日常生活方式，等等。这些变化在你侨居的国家或许也在发生，但那是在你在场时发生的，所以你能边走边适应。然而，那些并不是主要困难。真正的困难在于更小的日常琐事——你可能经历所谓“文化冲击”，虽算不上冲击，却同样令人困惑。比如超市结账时的确切做法可能变了；日常对话中参照框架也变了，你常常发现自己根本不知道别人在说什么，尽管他们说的是你的母语。}
\ptesub{题目 Question}
\begin{notebox}\small Which of the following are mentioned as being difficult to re-adapt to when returning to your native country after a long absence?\end{notebox}
\ptesub{选项 Choices（多选）}
{\small
A）Technological change\\
B）More than nine years absence\\
C）The difference in time zones\\
{\color{cFix}\textbf{D）The way people go about their daily lives}}\\
{\color{cFix}\textbf{E）People's terms of reference in conversation}}\\
F）Changes in government}
\smallskip\par
{\small 正确答案 \textcolor{cFix}{\bfseries D, E}\qquad 你的答案 {\color{cWrong}\bfseries A,} {\color{cFix}\bfseries D, E}\;{\footnotesize\color{gray}（D、E 对，多选 A 错扣分）}}
\end{qblock}

% =================== Q3 ===================
\begin{qblock}{Q3\quad FIB-L \#155\ \ Transport Chaos}{听力填空 Fill in the Blanks · 评分 \textcolor{cAccent}{\bfseries 1 / 4}}
\ptesub{录音原文 Transcript（含填空作答）}
\begin{promptbox}
The New South Wales government has \rok{apologized} for yesterday's transport chaos in and around Sydney Harbour during the visit of the Queen Mary II and the Queen Elizabeth II. Roads were jammed, traffic ground to a halt while tram and \rno{serious}{ferry} services were swamped with thousands of additional passengers, with most services delayed for hours. Premier Maurice Humor says that plans were put in place to ``deal with the congestion'' but the number of visitors well exceeded \rno{expantation}{expectation}. On the harbour itself there seemed to be as much congestion as there was on the roads, but everyone agreed it was an amazing \rno{spectical}{spectacle}.
\end{promptbox}
\zh{新南威尔士州政府就昨天玛丽女王二号与伊丽莎白女王二号到访期间、悉尼港周边的交通混乱致歉。道路被堵死、交通陷入停滞，有轨电车与渡轮服务被成千上万的额外乘客挤爆，大多数班次延误数小时。州长 Maurice Humor 称已制定方案“应对拥堵”，但访客人数远超预期。港口本身的拥堵程度似乎与道路不相上下，但人人都认同这是一场惊人的盛况。}
\par{\footnotesize\color{gray}（FIB-L 为听写打字，无选项；你的作答：apologized / serious / expantation / spectical —— 后三处分别听错词或拼写错误。）}
\end{qblock}

% =================== Q4 ===================
% （待填充）
